\section{Empirical Scope and Implementation}

The current paper is a methodological Monte Carlo study rather than an
application to a proprietary auction data set. An empirical implementation
requires auction-level information that distinguishes the highest submitted
bid from the transaction price, records the bidder count and auction format,
and identifies whether a posted reserve censored participation or sale. Repeated
auctions must be comparable after conditioning on observed item and market
characteristics.

With those data, estimation proceeds within covariate cells or through
conditional nuisance learners. Each training fold fits the structural
distribution and the local score moments; held-out auctions enter only the
orthogonal score. Bidder counts either index separate observation maps or enter
the structural simulator. The fixed-bandwidth target should be reported first,
followed by the bandwidth path, root count, and bias-aware exact-reserve
interval. Counterfactual revenue calculations use the global structural fit,
whereas uncertainty for the reserve uses the local score learner.

The maintained identification assumptions exclude endogenous entry,
unobserved auction heterogeneity correlated with bidder values, affiliated
private values, and confusion between the winning bid and transaction price.
An application lacking information to address these issues would not validate
the baseline model merely by matching the marginal distribution of observed
prices.
